%%%%%%%%%%%%%%%%%%%%%%%%%%%%%%%%%%%%%%%%%%%%%%%%%%%%%%%%%%%%%%%%%%%%%%%%%%%%
% forced-objects-zh.tex  (v4：按可读性重写)
%
% 一段一个意思，先说人话，再上符号。形式陈述放进表格与"方法"，叙事承担论证。
%
% 诚实账本（与 v3 一致）：
%   * "已验证" = Lean 4 声明，无 sorry，除 O1 外无具名公理；
%     native_decide 的用法在"方法"里列出，它支撑的是本文不报告的结果。
%   * 一条具名公理 O1，且没有定理消费它。
%   * 前提 C0（语架必须是标准的）未形式化；需要它的陈述标"模 C0"。
%   * 新颖性测量被本文自己的例子推翻，按失败如实报告，不当作可用闸门。
%%%%%%%%%%%%%%%%%%%%%%%%%%%%%%%%%%%%%%%%%%%%%%%%%%%%%%%%%%%%%%%%%%%%%%%%%%%%

\documentclass[aps,prd,preprint,superscriptaddress,amsmath,amssymb]{revtex4-2}
\emergencystretch=2.5em
\usepackage{graphicx}
\usepackage{bm}
\usepackage{dcolumn}
\usepackage{amsthm}
\usepackage{booktabs}
\usepackage{color}
\usepackage[fontset=fandol]{ctex}

\newtheorem{theorem}{定理}[section]
\newtheorem{definition}[theorem]{定义}
\newtheorem{remark}[theorem]{注}
\newtheorem{obligation}[theorem]{义务}

\newcommand{\lean}[1]{\texttt{\small\detokenize{#1}}}

\begin{document}

\title{数学对象什么时候是"被迫"的？\\
一个检验、三种裁决，和一次失败的测量}

\author{Yuan-Jie Liu}
\email{yuanjieliu65@gmail.com}

\date{\today}

\begin{abstract}
数学家常说某个对象\emph{必须}存在：自然数里减不了，于是有整数；实数里解不出
$x^2+1=0$，于是有复数。同一批对象又常被叫作发明，而这两种说法通常都不带检验。
本文提出一个检验，形式是几个只能回答"是/否"的问题，答案都在证明助手里核过。
问题一：旧系统能不能在不改变"元素是什么"的前提下满足这个要求？
问题二：在所有扩张方式里，有没有一个最小的——小到任何别的扩张都唯一地由它映射过去？
问题三：相对一份已知对象的清单，它是不是新的？
问题二答"是"时，对象就被钉住了：任意两个最小扩张只差一次改名，而这次改名本身唯一。
我们拿这套检验去查自己的工作：数系塔有三步通过了问题一与问题二；从有理数到实数那一步
没有通过，我们把它缺的性质写成一条点名义务，而不是假装证了；纲领自己提出的一个计数
不变量直接卡在问题一，因为它已经能在同一批元素上定义出来。
问题三则是失败的地方：我们实现"是不是新的"时比的是元素个数，而整数、有理数、复数
元素一样多，于是检验把该收的全拒了。这个失败我们照实报告，因为它恰好说明检验有牙齿。
活下来的是两个能用的问题、一个没实现的第三问，以及一份"这套办法做不到什么"的清单。
\end{abstract}

\maketitle

\begin{quote}
\itshape
凡可被定义者，皆可被追问。
\par
\vspace{0.4em}
--- 《银河帝国》
\end{quote}

\section{问题}
\label{sec:question}

\paragraph{一个答案错了的例子。}
假设有个程序宣布它发现了一个新数：$\alpha$，满足 $\alpha^2=7$。这话没有错，
但它不是发现。有理数里确实没有 $7$ 的平方根，所以 $\alpha$ 确实扩张了有理数；
可这个扩张就是大家都知道的那个域，只是把 $\sqrt 7$ 写成了 $\alpha$。
一个只问"以前有没有解"的检验会把它判成新的，而且会把一大堆东西都判成新的。

数系的通常说法也有同样的毛病。"自然数解不了 $x+1=0$"是真的；单点集、自由群、
有理函数域也都解得了这个方程，可它们都不是整数。所以
\emph{以前没有解，于是我们添一个}这句话并没有挑出一个对象，它挑出了一群。

本文要补的就是这句话缺的那半截。我们想做一个检验：别人（更有用的是自己）提一个
"这是发现"的说法时，能照着问几个问题，得到的答案不依赖于对象是怎么造出来的。
检验做出来之后，我们拿它查了自己的纲领：它拒掉了我们一个构造，拒绝认证另一个，
在第三个问题上直接坏掉。三件事都写在这里。

\paragraph{一句话的想法。}
一个对象是\emph{被迫}的，当这个要求没法在原来的元素上满足，而满足它的那个扩张是
\emph{最小}的——最小不是指规模小，而是指：任何别的扩张都有唯一一个映射从它过来。

"最小"这个词要小心，因为多数"必然性"的说法就是在这里偷偷作弊的。
拿"$x+1=0$ 必须有解"这个要求来说。如果我们允许继续待在"加法只是半群运算"的结构里，
就可以添一个形式符号 $\varepsilon$ 让 $\varepsilon+1=0$ 然后收工。
它确实解了方程，它不是整数：它的加法不是群运算。反过来，如果我们先同意在交换环里做事，
同一个要求就产生整数，而且只产生整数。差别不在要求里，而在要求被允许待的那家公司里。

所以检验必须先固定"公司"。这是整件事里最有用的一条教训，也容易说：
\emph{换背景理论就换对象；那一步的真正内容，是换理论，而不是添元素。}

\subsection*{背景：同一个对象，很多种造法}

数系是这个问题最老的例子，而它之所以被反复拿来当例子，恰恰因为答案不在造法里。
整数可以造成"自然数对模一个差"，也可以造成加法幺半群的群完备化；有理数可以造成整数对，
也可以造成整数的分式域；实数可以用 Dedekind 切割（1872），也可以用 Cauchy 序列的等价类；
复数可以造成实数对（Hamilton，1837），也可以造成商环 $\mathbb{R}[x]/(x^2+1)$。
任意两种造法给出的对象同构。Landau 的《分析基础》（1930）从 Peano 公理出发把整座塔造起来，
并给出标准的唯一性陈述；Ebbinghaus 等人合著的《Numbers》（1991）选讲各系统的多种造法及其等价性，据我们所知，一位数学家想把整座塔放在一处看时，拿的就是这本书
\cite{landau,dedekind,ebb}。

这批文献解决的是\emph{对象是什么}。它没有解决数学家引介这些对象时真正说的那句话——
那句话是因果的，不是结构的：\emph{我们减不了，所以造了整数}。
有意思的正是这句因果句，因为被滥用的也是它。一个程序添了个形式平方根 $\sqrt 7$，
可以一字不差地说同一句话，而它是错的。

这个困难的哲学版本比形式版本更老。Benacerraf（1965）指出：如果数就是集合，
那自然数 4 既可以是 von Neumann 的 4，也可以是 Zermelo 的 4，而算术里没有任何东西告诉我们
该选哪个；于是要么这个等同是任意的，要么数根本不是对象 \cite{benacerraf1965}。
结构主义的回应是：数是结构里的位置，不是有内在本质的对象 \cite{shapiro1997}。
我们对这场争论不表态。我们从里面取出一条设计约束：无论提出什么检验，
对同一个对象的每一种造法都必须给出同一个答案。本文的第二问就是在执行这条约束，
也是为什么检验的答案是关于泛性质的陈述，而不是关于某个集合的陈述。

\subsection*{关于"必然性"，前人已经说过什么}

数学家对\emph{必须存在}其实有个标准说法：对象在相应范畴里是\emph{初始}的
（或\emph{自由}的、或\emph{泛}的）。Mac Lane 的教科书陈述并使用了"初始对象唯一到唯一同构"
这条事实 \cite{MacLane1971}。具体实例也都是教科书内容：整数是加法幺半群的群完备化，
有理数是分式域，复数是"由一个 $x^2+1$ 的根生成的环"。模型论用另一套语言给出同一形状：模型完全存在时唯一到等价，比较嵌入在基模型上初等且唯一 \cite{EklofSabbagh1971,Hodges1993}。

所以材料并不新，这一点我们写在正文里，而不是留给审稿人来说。就本文所引文献而言，缺的是一个\emph{检验}：
一串短问题，读者能拿去问任何"这个对象必须存在"的说法——包括程序提出的说法——
外加一份诚实交代：哪些问题能机械回答，哪些不能。教科书的陈述是对一个固定范畴说的；
检验必须能扛住换范畴，因为滥用恰好发生在换范畴的地方。

\subsection*{把问题说清楚}

\begin{table}[t]
\caption{本文处理的问题，按我们希望审稿人核对的形式写。最后一行区分三级证据，
而不是笼统声称。}
\label{tab:problem}
\begin{ruledtabular}
\begin{tabular}{p{0.14\linewidth}p{0.74\linewidth}}
项 & 陈述 \\
\hline
已知 & 一个结构 $S$（带运算的一批对象），以及一个 $S$ 满足不了的要求 $D$。\\
要问 & 是否存在一个被 $D$ \emph{迫使}的对象——任何可接受的答案都必须是它，至多差一次唯一改名？\\
障碍 1 & $D$ 可能在同一批对象上已经满足，那就什么都没被逼出来。\\
障碍 2 & 可能有大量扩张满足 $D$，包括平凡的与混沌的；没有"最小"的那个，"那个"答案就没有定义。\\
障碍 3 & 答案可能只是已有对象换了个名字。\\
我们的答法 & 三个问题，各对一个障碍：同一批对象上是否可解；是否存在唯一到改名的最小扩张；
相对一份显式清单是否算新。\\
实际交付 & 三级证据分开写：机器检验（塔的三步、四处失败、两个打分的分离）；
手工论证（对计数不变量的拒绝）；断言（$\mathbb{R}$ 的极小性）。问题三只给规格未实现，
且第一版实现被本文自己的例子推翻。\\
\end{tabular}
\end{ruledtabular}
\end{table}

这样一来，论文的问题就窄到可以回答了：给定一个结构与一个要求，
是否存在被这个要求迫使的对象？以及，我们能不能——用怀疑者可以重复的方式——
判断答案什么时候是"是"？

\subsection*{先给结论}

四条，写在这里，好让只读这一页的读者也知道发生了什么。

\begin{itemize}
\item \emph{塔有三步通过。} 整数、有理数、$(\mathbb{C},i)$，在背景理论固定之后，
      各自是前一个系统的最小扩张，并且都唯一到改名。这些都是经典事实；
      贡献在于同一句话覆盖了这三个。
\item \emph{有一步无法认证。} 实数未被证明是有理数之上最小的完备有序域。
      我们证了引发这一步的失败，也证了刻画实数的切割性质；缺的那句话登记为一条点名义务，
      而开发里没有任何东西依赖它。
\item \emph{有一个候选被拒。} 纲领自己提出的计数不变量塌缩成一个素因子计数差，
      因而在同一批对象上已经可定义。检验拒绝了它，方式是手工论证。
\item \emph{第三个问题失败了。} 我们实现"是不是新的"时比的是元素个数，
      这使整数、有理数、复数互相成为对方的改名；闸门于是拒掉一切，
      包括本文其余部分刚刚认证的对象。
\end{itemize}

\section{前人做了什么}
\label{sec:related}

有四批工作和这个问题相关，本文与它们的分界值得明说。

\paragraph{造出数系，并证明结果唯一。}
这些造法是十九世纪到二十世纪中叶的数学：Dedekind 切割、Cantor 序列、Hamilton 的数对，
以及把半环变成环、把环变成域的代数构造。系统整理并给出唯一性证明的是 \cite{ebb}；
从 Peano 公理出发的公理化处理是 \cite{landau}。这批文献是本文每个例子的底账，
关于对象本身我们不声称任何新东西。

\paragraph{泛性质与比较态射的唯一性。}
一般陈述——初始对象唯一到唯一同构——见 \cite{MacLane1971}，而且到处在用：
自由群、局部化、多项式环、商。模型论用另一套语言给出同一形状，
经由模型完全与比较映射的唯一性 \cite{EklofSabbagh1971,Hodges1993}。
我们的主定理是这条陈述的参数化版本：背景理论与强化都是参数，
所以同一个四行论证覆盖整座塔。正文里我们明说这是打包，不是新数学。

\paragraph{形式化。}
实数在证明助手里被造过——HOL Light 是一例 \cite{Harrison1998}，更早还有 HOL88；这里用的代数层级是 Lean 证明器 \cite{lean} 上的 Mathlib \cite{mathlib}，本文的开发用 Lean~4。让这种层级可用的实例参数经验见
\cite{Baanen2022}。证明助手给这类论文加的不是深度，而是\emph{可审计性}：
定义没法悄悄变成另一个意思；而"这个证明只有四行"这种说法，可以当成关于文件的事实来核，
不必凭信。我们的体会是，这份代价主要买到了一件事——让错误显形。
本文里就有我们自己的两处。

\paragraph{机器发现。}
提议数学的程序有很长的历史，从 Lenat 的 \textsc{AM} \cite{lenat1976}，
到近年用搜索或学习生成猜想与构造的系统 \cite{davies2021,raayoni2021,romera2024}。
它们的样子很一致：先生成候选，再决定哪些值得留。这个"决定"由人来做，或由"有趣性"启发式来做（如 \textsc{AM}），
或由一个只验真、验性能、不验新的检查器来做。本文管的正是这个决定步骤。
它不生成任何东西，应当与那些系统互补着读：生成器提供候选，判据决定哪些可以叫被迫。

\paragraph{哲学：发现、发明，以及"数不能是什么"。}
本文这个问题所属的那场争论是 Benacerraf 的 \cite{benacerraf1965}，
结构主义的回应见 \cite{shapiro1997}。我们不去了结它。我们只取一条：
若答案要与造法无关，检验就必须在换造法之下不变。这是一条形式要求，
执行它的是本文的第二问。"被迫"的哲学读法仍然只是一种读法；
我们能核的是对象被确定了。

\paragraph{本文站在哪里。}
不在对象的构造里，不在唯一性定理里，也不在形式化里。它站在组合里：
一个短检验，写成对任意理论都适用的样子；诚实地用在作者自己的纲领上；
否定结果与肯定结果用同样的篇幅报告。最近的邻居是范畴论里的初始性陈述；
差别在于我们追问的是"初始性声称\emph{为假或未证}时会发生什么"——
而那恰好是一位在做数学的数学家通常所处的处境。


\section{一个含三问的检验}
\label{sec:test}

问任何问题之前，先固定三样东西：一个理论（新对象允许携带什么）、一种把旧结构读进
新结构的方式、以及一个要求（要添的常元、要加的性质，或两者都有）。固定之后，检验是三个问题。

\paragraph{问题一：这个要求是真的没被满足吗？}
问的是元素，不是解。这个措辞在干活。自然数与整数都是可数的，所以整数上的任何环结构都能沿一个双射搬到自然数上。
"自然数上不存在环结构"这句话因此是\emph{假}的，而我们自己的形式化早期版本就写了这句假话。
真的说法要把旧结构固定住：自然数上没有任何环结构，其加法乘法延续通常的加法乘法，
能解 $x+1=0$。这段插曲值得写下来，因为检验抓到的第一个人是作者自己。

\paragraph{问题二：这个扩张是最小的吗？}
在所有满足要求的扩张里，有没有一个能映射进任何别的扩张，而且这样的映射恰好只有一个？
答"是"时我们称这个扩张\emph{极小}。这个名字是传统叫法，性质其实不是"小"，而是"泛"；
而正是"泛"让对象被确定。下一节说清它带来什么。

\paragraph{问题三：它是新的吗？}
新，是相对已知的东西而言，不是相对系统自己的状态——系统自己的状态总能被安排成
"我这儿是新的"。实践中这意味着要有一份外部给的已知对象清单。
第~\ref{sec:gate}~节报告我们试着实现第三问时出了什么事；简短版是：实现错了，
而且只有拿例子去撞才发现。

\begin{table}[t]
\caption{三个问题、各自挡掉什么、以及本文实际交付了什么。"已答"指对本文的例子有机器
检验的结论；"只给规格"指我们写清了答案该是什么，但没有实现它。}
\label{tab:three}
\begin{ruledtabular}
\begin{tabular}{p{0.05\linewidth}p{0.32\linewidth}p{0.32\linewidth}p{0.17\linewidth}}
\# & 问题 & 挡掉什么 & 本文状态 \\
\hline
1 & 旧元素撑得起这个要求吗？ & 换掉元素却把它叫作解 & 已答 \\
2 & 有没有最小扩张，且唯一到改名？ & 平凡扩张、混沌扩张，以及一切依赖选择的构造 & 已答 \\
3 & 相对已知清单，它新吗？ & 给已有对象改个名字 & 只给规格，未实现；第一版实现失败（\S\ref{sec:gate}） \\
\end{tabular}
\end{ruledtabular}
\end{table}

\subsection*{这里的"被迫"是什么意思，不是什么}

这个词至少背着四种不同的主张，不说清是哪一种，就会四面挨打。按强度排：

\begin{itemize}
\item \emph{形而上学}：这个对象不可能不存在。我们不碰这一种。
\item \emph{概念}：对象已经包含在我们出发的概念里。这一种也不碰；我们的要求是数据，不是概念。
\item \emph{实践}：以本领域的现状看，没有能用的替代品。这是关于学术共同体的判断，不是关于数学的。
\item \emph{被确定}：任何可接受的答案都是同一个对象，至多差一次唯一改名。这是我们检验的那一种。
\end{itemize}

第四种比前三种弱，也是唯一能靠读证明来了结的。它同样是 Benacerraf 的论证留给结构主义的那一条：
如果数是结构里的位置，那么"同一个位置"就是唯一改名下的确定，再强的东西也拿不出来。
所以我们不声称证了整数必然存在。我们声称的是：给定这个要求和它被允许待的公司，
没有别的东西能应这个要求。

这不是留后路，它改变了什么算反驳。想要形而上学必然性的读者会觉得这套检验沉默，
它也确实该沉默；想知道一个程序的"发现"是不是改名的读者，会觉得它说了算。

\section{为什么第二问能把对象钉住}
\label{sec:theorem}

问题二有一个干净的答案，而这个答案正是这套检验值得做的原因。如果一个扩张在上述意义下
最小，那么：

\begin{enumerate}
\item 它映射进任何别的扩张，且这个映射唯一；
\item 任意两个最小扩张，只差一次对其元素的改名，且这次改名本身唯一；
\item 最小扩张自己没有恒等以外的对称。
\end{enumerate}

这三条没有一条提到环、序或完备性。它们只用到周围范畴的四条性质——恒等映射、复合、
把映射看成底层函数、以及"取值相同的映射相等"——别的什么都没用。这就是为什么一个论证
同时覆盖整数、有理数和复数，也是为什么它只需核一次，而不是三次。

它值多少，应该说明白。一旦你把"最小" \emph{定义}成"唯一地映射进一切"，
第一条和第三条就是定义的重述，证明只有四行。内容不在证明里。内容在于：
这个定义是当时桌上唯一不依赖构造方式的定义。两个人用不同办法造出复数——实数对、
商环、矩阵代数——对问题二会得到同一个答案，而他们的构造之间只差一次唯一改名。
这就是"被确定的是对象、不是构造"的意思。

\paragraph{一个例子，以及一个警告。}
从实数到复数那一步，要求是"添一个常元 $j$，满足 $j\cdot j=-1$"。
扩张 $(\mathbb{C},i)$ 是极小的，而且除恒等外没有对称。现在把常元从要求里拿掉，
只要求"存在某个 $-1$ 的平方根"。同一个域仍是扩张，但它不再极小：比较映射不再唯一。
恒等与复共轭是两个不同的 $\mathbb{R}$-映射，而且都被允许；把 $i$ 钉住，正是去掉第二个的原因。\emph{找到一个解}和\emph{把解钉住}是两件不同的成就，
只有第二件产生对象。一个接受未钉住版本的检验，会接受一个带对称群的域——
而那正是它本该消除的歧义。

\paragraph{这个论证不管的一件事。}
"序是完备的"这个要求和"$j^2=-1$"表现不一样，这一点用人话说更清楚。
像 $j^2=-1$ 这样的等式，任何环映射都保持它：$j$ 满足，$j$ 的像也满足。
像"每个有上界的非空集合都有上确界"这样的条件，环映射不保持它，而且它没有项模型——
你没法写下一串表达式、作个商、就完事。由此得到的是：最小扩张\emph{算不出来}，
不能像复数那样算。它\emph{不}意味着最小扩张不存在。我们的初稿下了更强的结论；
审稿人是对的，我们撤回了它（第~\ref{sec:real}~节）。

\section{检验查到了什么}
\label{sec:found}

我们把检验用在数系塔的四步上，以及自己纲领的另外两块上：一个由素因子分解造的不变量，
一个由可逆演化造的坐标。表~\ref{tab:results}~是汇总；标"不适用"的两行就是这两块，
理由是不是"对结构提的要求"——一个是定义，一个是坐标系——问题一、二没有咬的地方。

\begin{table}[t]
\caption{通过了什么、没通过什么、以及什么没法判。"同元素"是问题一，"最小"是问题二，
"新"是问题三。横线表示该问题不适用或没用到。}
\label{tab:results}
\begin{ruledtabular}
\begin{tabular}{lccc}
情形 & 同元素 & 最小 & 新 \\
\hline
$\mathbb{N}\to\mathbb{Z}$ & 已证 & 已证 & 未用 \\
$\mathbb{Z}\to\mathbb{Q}$ & 已证 & 已证 & 未用 \\
$\mathbb{Q}\to\mathbb{R}$ & 已证 & \emph{只是断言} & 未用 \\
$\mathbb{R}\to\mathbb{C}$ & 已证 & 已证 & 未用 \\
幂与开方 & 不适用 & 不适用 & 不适用 \\
计数不变量 & \emph{不过}（手工拒绝） & --- & --- \\
振动坐标 & 不适用 & 不适用 & 不适用 \\
\end{tabular}
\end{ruledtabular}
\end{table}

\subsection{通过的三步}
\label{sec:tower}

一旦同意在交换环里做事，整数就是自然数的最小扩张；一旦同意在域里做事，
有理数就是整数的最小扩张；一旦把 $-1$ 的平方根钉住，$(\mathbb{C},i)$ 就是实数的最小扩张。
底层事实是经典的（参~\cite{ebb}）；把它们写成"带旧结构钉住的最小扩张"是我们的包装。检验加的是：它们是同一句话的三个实例，而且每个答案都不依赖构造方式。
对整数那一步，还有一个附赠品：要求"扩张真的包含旧系统"会自动给出特征零，
这算是一个信号，说明定义在干活。

\subsection{一步我们没能认证}
\label{sec:real}

从有理数到实数这一步，我们证了失败——有理数不完备——也证了每个实数都由其下方的
有理切割生成。我们\emph{没有}证实数是有理数之上最小的完备有序域。
这句话在形式化开发里被写成一条点名义务，而开发里没有任何定理用到它。

我们认为这条义务为真、也可证；用标准的一页论证应该就够。我们没做。
后果是第四步的认证强度弱于前三步，而我们宁愿在正文里说清，也不愿让读者在附录里发现。
这个不对称是证明债与可计算性的不对称，不是"添根"与"完备化"在形而上学上分属两类。

\subsection{一个被拒掉的候选}
\label{sec:counting}

纲领还提出过一个整数的不变量，用来刻画两个数之间的"张力"：数一数只整除其中一个的素数，
减去反方向的个数，把结果叫 $\tau(a,b)$。提议是：曲面 $\tau=0$ 组织了算术，
而 $\tau$ 本身是一个新对象。

公理那部分做得很干净。在四条自然假设下，这个不变量唯一到符号，并且等于一个计数差：
$\tau(a,b)=\omega(a)-\omega(b)$，其中 $\omega(n)$ 是 $n$ 的相异素因子个数。
这就是它的全部内容：$\tau$ 知道两边各有几个素因子，永远不知道是哪几个。
由于 $\omega$ 是上百年的老函数、公式又能在同一批元素上定义出来，检验在问题一处
判\emph{拒}：这里没有新对象，只有一条关于 $\omega$ 的定理，
外加一句"这个构造为什么到此为止"的清楚说明。

两点如实交代。第一，这次拒绝是手工论证的：开发里计数那一层没有需求、也没有清单，
所以"问题一判拒"是数学观察，不是机器检验的判决。第二，这次拒绝不经过问题三；
初稿说这里的新颖性为零，那既推不出来，也不可得，因为清单根本不存在。

这段插曲仍然有用。它是检验对我们自己说"不"的那一次，而且带了理由：
不变量塌缩成一个计数，而计数已经在同一批元素上了。换一个更松的检验，
这里就会产出一篇声称发现的论文。

\section{失败的那次测量}
\label{sec:gate}

问题三是我们没能交付的那一问，而失败本身值得写。

设计里有两个分数。一个是五项（覆盖、关联、闭包增益、极小性、新颖）的加权平均，
用来\emph{排序}；另一个是四项（新颖、一致、生成力、压缩）的乘积，用来\emph{准入}。
区别是真的：加权平均会给"新颖性为零、其余都很好"的候选打高分，而那正是 $\sqrt 7$ 的故事；
乘积会给它零分。这条算术区别我们在证明助手里核过：存在新颖性为零、生成力 $4/5$ 的评分，
也存在同取值的快照其乘积为零（\lean{weighted_sum_admits_renaming}、
\lean{product_rejects_renaming}）。要紧的是说清它证明了什么、没证明什么。
它证明两个公式不可互换。它完全没给"改名"建模：那两条陈述里没有清单、没有候选、
没有对象之间的映射。

给改名建模，正是实现失败的地方。为了判断一个候选是不是新的，代码比的是载体的
\emph{元素个数}：如果候选的元素个数与某个已知对象相同，就叫改名。整数、有理数、复数
都是可数无穷。于是清单里只要有一个可数对象，这三者全成了它的改名，
闸门给三者的新颖性全打零分。闸门拒掉的，恰是本文其余部分认证的对象。

我们没有打补丁。比元素个数只会产生"错拒"，不会产生"错收"，所以风险方向与初稿写的相反。
正确的比较是在结构之间，而不是在载体之间，这要求清单自带映射；那是规格，我们没实现。
因此本文没有任何准入决定用到问题三，随代码发布的轨迹里印出的新颖性数值是手填的。
这句话放在表注里，表注也确实写了。

\section{这套检验做不到什么}
\label{sec:limits}

\paragraph{它假定映射是标准的。}
主论证用的四条性质没有说结构之间的映射是"保持结构"的那些。一个退化的选择——
任意两个结构之间只有一个映射——也满足这四条，而在那种情形下每个对象既最小又最大，
检验对任何事都不发言。我们的例子按构造使用标准映射，但这条假定在框架内部证不出来，
只能逐实例核对。凡需要它的陈述，我们标"模 C0"。这是把判据写成"任意理论都适用"
所付的代价。

\paragraph{它不生成需求。}
发现引擎里难的部分——判定需求是否可满足、提出候选、判定极小性、判定一致性——
在本文里是输入，不是代码。我们证的是它们周围的事：可解的需求不需要新元素；
矛盾永远不会被扩张；旧系统嵌在新系统里；同一需求的两个发现只差唯一一次改名。
引擎是法官，不是作者。

\paragraph{它不能判定自己的问题。}
每个问题都是关于一组固定数据的命题。判定它是逐实例的活，而对问题一来说，
这活就是"在一个结构里解方程"。本文没有任何判定过程，也不声称三个问题一般可以机械回答。

\paragraph{它只拒过我们一个候选。}
另一个情形，实数那一步，是"未认证的义务"而不是"拒绝"。只说过一次"不"的检验是开头，
不是战绩。我们最想拿出的证据，是拒掉别人提出、并且他本人相信的候选；这个现在还没有。

\paragraph{这套检验是相对于数据的。}
本文每个裁决都相对于三个选择：背景理论、强化、以及（对问题三）清单。改动任意一个，
裁决都可能变。这不是待修的缺陷，这正是裁决可核查的原因。但它确实限定了主张：
我们没有证明某个对象\emph{无条件地}是被迫的，只证明了它相对于一组说清的三元组是被迫的。
想要无条件那个词的读者在本文里找不到，而诚实的回答是：我们不知道那个词该是什么意思。

\section{讨论}
\label{sec:discussion}

实践上的教训是关于解释的顺序。通常的顺序是：先给你看对象，再给你看使它必要的性质。
我们的检验活下来之后，顺序反过来了。先给一个在原地修不好的失败。再给"允许在什么公司里修"，
因为同一个失败在不同公司里修出来的对象不同。再给最小的那个修法——公司一固定，它就定了。
这三步之后剩下的东西不是被迫的，应该叫构造。

第二条教训是：否定结果很便宜，而且值得发表。纲领的计数不变量曾经是另一篇论文的主角。
拿它过一遍问题一花了一个下午，产出的是关于 $\omega$ 的定理。这比一篇声称发现的论文
结果更好，成本更低。

第三条教训关于机器检验。证明助手没有解决任何哲学问题，也没有让数学变深。
它做了三件有用的事：逼我们把"缺口"的定义写得不是假的；让主论证只有四行这件事显形，
使我们没法假装它更厚；把"两个打分公式不可互换"这句声称变成一条一分钟能核完的定理。
它也让我们的错误显形——环结构那处、比元素个数那处——这才是这类论文里用它的主要理由。

\paragraph{什么会让我们改口。}
三件事。第一，一个真正的发现，通过了问题一与问题二，却被修好的问题三拒掉——
那说明第三问错了，而不只是没实现。第二，某个需求的扩张不构成范畴——
那说明第二问一般不可用。第三，有人证出实数不是有理数之上最小的完备有序域——
那让表里第四行从"未证"变成"假"。第三件我们预计会往反方向走：那条义务看起来可证，
而我们会乐于撤回这个不对称。

\paragraph{与已有工作的关系。}
问题二的唯一性那一半是范畴论常识~\cite{MacLane1971}；它的模型论表亲——
模型完全给出的比较唯一性——见~\cite{EklofSabbagh1971}。本文加的是打包方式：
对任意理论、把需求与背景理论当参数的同一个陈述，以及由它得出的两个否定裁决。
自动猜想系统~\cite{lenat1976,davies2021,raayoni2021,romera2024}生成命题，
由命题是否为真、是否有用来评判；它们提供候选，本文管另一头：决定哪些候选可以叫被迫。

\section{方法}
\label{sec:methods}

\paragraph{形式设置。}
一个理论把每个类型送到"该类型上的结构"的类型；一个模型是载体加结构。一个语架提供
模型之间的映射、恒等、复合、底层函数映射，以及外延性（取值相同的映射相等）。
复合的结合律与单位律由这些推出，不作假设。一个需求为每个模型提供见证的类型、
见证上的谓词、以及沿映射搬运见证且保持函子性的操作。模型的一个扩张是：目标模型、
从旧载体出发的单射、以及满足谓词的见证。它\emph{极小}，当它唯一地映射进任何别的扩张。
第~\ref{sec:theorem}~节的三条结论就是开发中的定理一（\lean{IsMinimal.hom_unique}、
\lean{IsMinimal.carrierEquiv}、\lean{IsMinimal.rigid}）。

\paragraph{信任基底。}
开发是 Lean~4~\cite{lean} + Mathlib~\cite{mathlib}，无 \texttt{sorry}。有一条具名公理，
即实数的初始性，没有任何定理消费它。另外，\lean{native_decide} 出现在计数层的两个文件里；
它把 \lean{Lean.ofReduceBool} 引入信任片段，支撑的是本文不报告的有限数值核对。
本文从那一层真正引用的陈述，是靠改写与算术证的，不是靠求值；初稿把这件事说反了，
更正记录在仓库里。

\paragraph{哪些核过、哪些只是论证。}
核过（每一条都有具名声明，名字收在审计文件里）：主定理；四处失败；三条极小性结果；
特征零的附赠条件；三对判定中两对的互斥引理；引擎的四条正向性质；幂映射的挠分解；
计数不变量塌缩成计数；两个打分的分离。
论证而未核：等式被保持、闭包条件不被保持；对计数不变量的拒绝；缺口分支里候选与旧世界的比较。
假设：标准映射前提（C0），以及实数的初始性。

\paragraph{数据与代码。}
Lean 开发与本文的七轮自审都在仓库里，打了 tag \lean{paper-forced-objects-v1}；
构建把 Mathlib 钉在 \lean{lakefile.pinned.toml} 里；文中点名的每个声明都能按标识符定位，
审计记录写明哪些声称被撤回或降级、以及为什么。有三条边界要照直说。
计数层与振动层没有单独的可执行模型：它们是 Lean 定义；\lean{Nods/Chain.lean} 里那份轨迹是
Lean 文件，分数是手填的，不是程序算出来的。仓库里还放着与本文无关的工作（一个模拟研究），
而且 Mathlib 的版本没有钉死，所以我们不宣称"一条命令就能构建"的 artifact。
文中点名的每个声明都能按标识符定位；审计记录写明哪些声称被撤回或降级，以及为什么。

\begin{acknowledgments}
开发在 Lean~4 + Mathlib 中完成，感谢 Mathlib 社区。本文经七轮对抗式自审，由作者的工具以独立审稿线的形式跑出；
审计记录（含两处事实性撤回与第三问的降级）与源码保存在一起。
\end{acknowledgments}

\section*{利益冲突}
作者声明无利益冲突。

\bibliography{nods}

\end{document}
